\section{Problemstellung}
\label{sec:problemstellung}

\subsection{Fachliche Betrachtung}
\label{subsec:fachliche-betrachtung}

Beim Rollenspiel \Abk{DSA} handelt es sich um ein \sog \Abk{PPRPG}.
Diese zeichnen sich \insb dadurch aus, dass die Teilnehmer physisch am selben Ort sind und vor allem mit Stift und
Papier gearbeitet wird.
Grundsätzlich wird in einem \Abk{PPRPG} eine fiktive Geschichte von fiktiven Charakteren durch die Teilnehmer interaktiv
konstruiert und erzählt.
Ein Teilnehmer besetzt die Rolle des \sog Spielleiters\; die Restlichen sind \sog Spieler.
Der Spielleiter führt anhand einer fertig beschriebenen Spielwelt und einem passenden Regelwerk die Spieler durch die
Geschichte.
Die Spieler besetzen jeweils die Rolle eines Charakters und beschreiben diesen auf einen \sog Charakterbogen, auf
dem \uA dessen Beruf und Fähigkeiten stehen.

Das \Abk{PPRPG} Regelwerk zu \Abk{DSA} ist im Jahr 1984 erschienen und befindet sich aktuell in der 5ten-Ausgabe aus
2015.
Es zeichnet sich \insb durch den Umfang und die Komplexität aus, was \uA die Character-Erstellung und -Verwaltung für
die Teilnehmer deutlich schwer macht.

\subsubsection{Charakter-Erstellung \textit{(vereinfachte Darstellung)}}
\label{subsubsec:charakter-erstellung}

Ein Character definiert sich durch die folgenden Merkmale: seine Spezies, Kultur, Profession, Eigenschaften,
Fertigkeiten sowie Vor- und Nachteilen.
Einige Merkmale können Stufen \bzw eine Wertigkeit haben, \uA Eigenschaften, Fertigkeiten sowie Vor- und Nachteile.
Fertigkeiten bestehen aus Talenten, Zaubern, magischen Handlungen und Liturgien.
Jeder dieser Merkmale, oder die Steigerung einer Stufe/Wertigkeit, kostet dem Spieler \Abk{AP} und kann bei der Wahl von
Nachteilen \Abk{AP} zurückbekommen.
Bei der Character-Erstellung hat jeder Spieler eine, vom Spielleiter, vorgegebene Anzahl an Start-\Abk{AP}.
Zunächst wird damit die Spezies, Kultur(en) und Profession ausgewählt.
Mit den restlichen \Abk{AP} wird für Fertigkeiten, Vorteile und andere Spielmechaniken bezahlt oder durch Nachteile
erworben.
Abschließen sind acht Basiswerte zu berechnen.
Diese werden durch die Spezies, den Eigenschafts-Stufen, Vor- und Nachteilen sowie Sonderfertigkeiten beeinflusst.
\ZitatVgl{Regelwiki-Heldenerschaffung}

Dabei ist aus folgendem Umfang zu wählen (dies stellt ein Auszug und keine komplette Aufstellung dar)
\ZitatVgl[ff.]{Regelwiki}:
\begin{multicols}{2}
    \begin{itemize}
        \item vier Spezien
        \item 32 Kulturen
        \item 51 Professionen
        \item acht Eigenschaften
        \item 59 Talente
        \item 251 Zaubern und 74 Ritualen
        \item 181 magischen Handlungen
        \item 74 Liturgien
        \item 130 Vorteile
        \item 77 Nachteile
        \item weiteren Spielmechaniken wie \zB 21 Kampftechniken und 354 Sonderfertigkeiten
    \end{itemize}
\end{multicols}

Die Komplexität entsteht vor Allem durch die Vielzahl der Merkmals-Kombinationen und Regeln.
Die \og Merkmale können sich \bspw gegenseitig beeinflussen oder ausschließen:
\begin{itemize}
    \item Eine Spezies gibt Boni und Mali für Eigenschaften vor \ZitatVgl{Regelwiki-Spezies-Elfen}.
    \item Eine Spezies kann Vor- und Nachteile vorgeben \ZitatVgl{Regelwiki-Spezies-Elfen}.
    \item Eine Profession hat Voraussetzungen \uA \ggfs an Eigenschaften, Fertigkeiten, Vor- und Nachteilen sowie
    Sonderfertigkeiten \ZitatVgl{Regelwiki-Beruf-BruderDesFeuers}.
    \item Eine Profession gibt Fertigkeiten (Talente, Zauber, Rituale, Liturgien) vor
    \ZitatVgl{Regelwiki-Beruf-BruderDesFeuers}.
    \item Einige Fertigkeiten stehen nur bestimmten Professionen zur Verfügung (\sog Berufsgeheimnis*se)
    \ZitatVgl{Regelwiki-Berufsgeheimnis}.
    \item Eine Sonderfertigkeit hat Voraussetzungen \uA \ggfs an Eigenschaften, Fertigkeiten, Vor- und Nachteilen sowie
    Sonderfertigkeiten \ZitatVgl{Regelwiki-Sonderfertigkeit-Blutrune}.
\end{itemize}

Weitere Regeln, welche bei der Charakter-Erstellung zu beachtet sind, werden nachfolgend (und nicht
abschließend) aufgezeigt:
\begin{itemize}
    \item Es darf eine vorher festgelegte maximale Anzahl an aktivierten Zauber und Liturgien, sowie ``Fremdzauber``
    (nicht der Profession zugeordnet) nicht überschritten werden \ZitatVgl[Schritt 2]{Regelwiki-Heldenerschaffung}.
    \item Eigenschaften, Fertigkeiten und Kampftechniken können nur bis zu einem vorab festgelegten maximalen Wert
    verbessert werden \ZitatVgl[Schritt 2]{Regelwiki-Heldenerschaffung}.
    \item Die erste Kultur ist gratis \ZitatVgl[Schritt 4]{Regelwiki-Heldenerschaffung}.
    \item Zauber und Liturgien sind vor der Wert-Steigerung zu aktivieren, dafür sind Aktivierungskosten zu zahlen.
    \ZitatVgl[Schritt 8]{Regelwiki-Heldenerschaffung}
    \item Es darf nur bis zu einer bestimmten Anzahl \Abk{AP} Vorteile gekauft oder mit Nachteilen erworben werden
    \ZitatVgl{Regelwiki-VorUndNachteile}.
    \item Am Ende darf nur eine maximale Anzahl \Abk{AP} übrig bleiben \ZitatVgl[Schritt
    11]{Regelwiki-Heldenerschaffung}.
    \item Eine Spezies gibt übliche und unübliche Kulturen sowie Vor- und Nachteile vor.
    \ZitatVgl{Regelwiki-Spezies-Elfen}
    \item Eine Profession gibt übliche und unübliche Kulturen sowie Vor- und Nachteile vor.
    \ZitatVgl{Regelwiki-Beruf-BruderDesFeuers}
\end{itemize}

Bei den letzteren Zwei handelt es sich um Vorgaben der Spielwelt, damit die Merkmals-Kombinationen in dieser stimmig
sind.
Unübliche Merkmals-Kombinationen sind nicht per se verboten, müssen jedoch bekannt und mit dem Spielleiter abgesprochen
werden.

\subsubsection{Charakter-Verwaltung \textit{(vereinfachte Darstellung)}}
\label{subsubsec:charakter-verwaltung}

Im laufe der Geschichte gewinnt man AP und kann sich damit steigern

\subsection{Technische Betrachtung}
\label{subsec:technische-betrachtung}

Bereits im Jahr 1993 wurde vom damaligen Herausgeber eine Software mit dem Namen ``\Abk{DSA}-Tools`` zur Erstellung
und Verwaltung von Charakteren für die damalige Ausgabe veröffentlicht und 1995 in einer neuen Version
\Zitat{Wiki-DSA-Tools,Wiki-DSA-Tools-Deluxe}.
Doch die aktuell genutzten Softwares kommen aus der Community selbst, nachfolgend die populärsten aufgelistet:
\begin{itemize}
    \item \textbf{Optolith Heldengenerator}, erschienen am \Datum{2018-04-08} und zuletzt Aktualisiert am
    \Datum{2022-03-22}, dient der Charakter-Erstellung und -Verwaltung in der 5ten-\Abk{DSA}-Ausgabe.
    Es ist seit \Monat{2020-09} OpenSource und wird hauptsächlich von einer Person gewartet.
    Genutzt werden die Sprachen TypeScript\Fussnote{Typisierter Superset von JavaScript; https://www.typescriptlang
    \item .org/} und ReasonML\Fussnote{OCaml (https://ocaml.org/) in JavaScript-angelehnter Notation;
    https://reasonml.github.io/}.
    Die Regeln werden in \acs{YAML}-Dateien festgehalten und ein Zugriff wird, aufgrund von Lizenzgebungen, nur auf
    Anfrage und Notwendigkeit ermöglicht.
    \Zitat{GitHub-Optolith}
    \item \textbf{Helden-Software}, erschienen im Jahr \Jahr{2012} und zuletzt Aktualisiert am
    \Datum{2017-07-23}, dient der Charakter-Erstellung und -Verwaltung in der 4ten-\Abk{DSA}-Ausgabe.
    Es ist nicht OpenSource.
    Genutzt wird die Sprache Java\Fussnote{Betriebssystem-unabhängige \Abk{OOP} Sprache; https://www.oracle
    .com/de/java/} und keine weiteren Frameworks/Bibliotheken \Zitat{HeldenSoftware-QA-Entwicklungsumgebung}.

    Es wird von einem Team aus sieben Personen (fünf Entwickler und zwei Tester) gewartet \Zitat{HeldenSoftware-Team}.
    Als Besonderheit bietet diese Lösung das Speichern der Charakter auf deren Server und eine Schnittstelle zur
    Entwicklung von Plugins an, mit denen Informationen zu den Charakteren abgerufen werden können
    \Zitat{HeldenSoftware-Internet,HeldenSoftware-Plugins}.
\end{itemize}


%%%  Community Wiki:
%%%  https://de.wiki-aventurica.de/wiki/Willkommen_im_Wiki_Aventurica!

%%%  Regelwerk vom Publisher
%%%  https://www.ulisses-regelwiki.de/start.html
%%%
%%%  Arten von Rollenspiele:
%%%  https://de.wikipedia.org/wiki/Rollenspiel_(Spiel)
%%%
%%%  Defenition PP:
%%%  https://de.wikipedia.org/wiki/Pen-%26-Paper-Rollenspiel
%%%
%%%  PP wissenschaftlich:
%%%  https://de.wikipedia.org/wiki/Rollenspieltheorie

